\subsection{Continuous-dynamics counterpart: the bistability needs near-full-strength wg self-activation}
Verdict \#4 asks whether the attractor conclusions survive continuous dynamics. We built a Hill-function ODE counterpart of the 5-node core (production: OR over activators, multiplicative repression; linear decay) and scanned $3^5 = 243$ initial conditions over 12 Hill-parameter settings ($K \in \{0.2, 0.35, 0.5, 0.65\}$, $n \in \{2, 4, 8\}$). \textbf{At every tested setting the ODE has a single all-OFF attractor: the wg-ON basin fraction is 0.} The mechanism is structural, not parametric: in the Boolean wg-ON state $(1,1,1,0,0)$, ci is OFF, and ci is wg's only activator in the core - so under continuous dynamics wg production vanishes and linear decay erases the state. The discrete bistability is update-rule-supported (consistent with the 14$\times$ async/sync accessibility gap above), not an intrinsic property of the wiring under mass-action-like dynamics.

The constructive version of this negative is a threshold: adding wg auto-activation of relative strength $a$, the wg-ON attractor reappears only for $a \geq 0.9$ of full Hill strength (basin 66.7\% of the IC grid at $a = 0.9$, versus 0 at $a = 0.85$ and below, over all 12 $(K, n)$ settings). Biologically this is a sharp, testable requirement: \textbf{if segment-polarity bistability operates continuously, wg must be near-fully self-sustaining} - a prediction consistent with the known wg autoregulation in the parasegment boundary and falsifiable by attenuating wg cis-autoregulation. Full record: \texttt{results/ode\_vs\_boolean.json}, \texttt{results/ode\_vs\_boolean\_scan.json}.
